\documentclass[tfg,estilo=digital,portada=normativa]{fdi-simplex}
% Cambiar arriba tfg por tfm en su caso

\titulo{Un sistema para ilustrar la plantilla fdi-simplex}
\tituloEN{A system to showcase the fdi-simplex template}
\autor{Frodo Bolsón}
% \calificacion{Notable (8)}           % Poner en la versión final
\autor{Samsagaz Gamyi}
% \calificacion{Sobresaliente (10)}    % Poner en la versión final
\director{Gandalf el Gris}
% \director{Elrond Medio-elfo}         % Se puede repetir
% \codirector{Elrond Medio-elfo}       % O también poner así
% \tutor{Arwen Tinuviel}               % En TFM tutor y cotutor

\titulacion{Ingeniería y esas cosas}
\cursoAcademico{2098/2099}

\palabrasClave{LaTeX, plantilla, TFG, Facultad de Informática}
\keywords{LaTeX, template, bachelor's thesis, Computer Science Faculty}

\addbibresource{bibliografia.bib}

% - OPCIONES DE PERSONALIZACION -

% Usar si se guarda el escudo en otra ubicación
%\logo{img/Escudo_UCM}

% Para cambiar las fuente tipográficas (solo LuaLaTeX/XeLaTeX)
%\setmainfont{Source Serif 4}
%\setmonofont{Source Code Pro}

\begin{document}

% Se puede cambiar el estilo de la portada con la opción
% portada=normativa|elegante en el documentclass, o se puede diseñar una propia
% (simplemente no llamar a este comando)
\makeportada

\frontmatter

\begin{resumen}
Este documento demuestra el uso de la plantilla \texttt{fdi-simplex}
para redactar la memoria de un Trabajo de Fin de Grado en la
Facultad de Informática de la Universidad Complutense de Madrid. La
plantilla incluye la portada normalizada, entornos bilingües para
resumen y palabras clave, e integración con biblatex.
\palabrasClave{LaTeX, plantilla, TFG, Facultad de Informática, Universidad Complutense de Madrid}
\end{resumen}

% En TFM, según la normativa punto 6 el resumen en inglés va primero, cambiar
\begin{abstract}
This document demonstrates the \texttt{fdi-simplex} template for
writing the report of a Bachelor's Thesis at the Faculty of
Computer Science of the Complutense University of Madrid. The
template provides a normalized cover page, bilingual abstract and
keyword environments, and biblatex integration.
\keywords{LaTeX, template, bachelor's thesis, Computer Science Faculty, Complutense University of Madrid}
\end{abstract}

\tableofcontents

\mainmatter

\chapter{Introducción}

Este capítulo introduce los antecedentes, objetivos y plan de
trabajo de la memoria, tal y como exige la normativa del trabajo
de fin de estudios de la Facultad de Informática
(\href{https://informatica.ucm.es/normativa-tfm-2425}{informatica.ucm.es/normativa-tfm-2425}).

\section{Normativa}

\begin{itemize}\setlength{\itemsep}{0pt}
  \item Extensión mínima: 25 páginas para TFG de un solo autor (más
        5 por cada coautor adicional); 50 páginas para TFM.
  \item Máximo 10 palabras clave en cada idioma.
  \item El resumen en inglés del TFM debe ocupar
        aproximadamente media página.
  \item Todo material no original (texto, figuras, código) debe
        citarse y respetar licencias.
  \item \textbf{Sólo TFM}: el estudiante debe firmar y adjuntar la
        \emph{declaración responsable sobre autoría y uso ético de
        herramientas de inteligencia artificial}.
\end{itemize}

\section{Antecedentes}

La Facultad de Informática de la UCM ha utilizado durante más de
una década el sistema TeXiS~\cite{texis} para sus memorias de
tesis y trabajos de fin de carrera. Esta cáscara ayuda a preparar
documentos complejos con buenas prácticas, pero al ser una
cáscara y no una plantilla, es más complicada de integrar con
otros sistemas, además de no incentivar el aprendizaje de LaTeX
por los usuarios.

\section{Objetivos}

Los objetivos de este trabajo son:
\begin{itemize}
  \item Disponer de una plantilla moderna basada en LuaLaTeX.
  \item Separar los aspectos normativos de los estéticos.
  \item Ofrecer una base mínima que no interfiera con los
        paquetes que cada estudiante decida usar.
\end{itemize}

\chapter{Desarrollo}

El cuerpo principal de la memoria describe el trabajo realizado.
En este ejemplo, demostramos algunas construcciones típicas: una
cita bibliográfica~\parencite{knuth1986} y referencias cruzadas a
la sección~\ref{sec:arquitectura}. La cita anterior es parentética
(autor y año dentro del paréntesis); cuando el autor ya forma parte
de la frase conviene usar \verb|\textcite|, que sólo encierra el
año: \textcite{tolkien1954lord} cuenta cómo Bilbo decía a menudo
que solo había un camino y que era como un río caudaloso; nacía
en el umbral de todas las puertas, y todos los senderos eran ríos
tributarios. "Es muy peligroso, Frodo, cruzar la puerta", solía
decirme. "Vas hacia el Camino, y si no cuidas tus pasos no sabes
hacia donde te arrastrarán".

\section{Arquitectura}
\label{sec:arquitectura}

La arquitectura del sistema sigue el patrón de separación entre
normativa y estilo. La clase \texttt{tfe-ucmfdi} implementa los
requisitos obligatorios de la normativa UCM-FDI; el aspecto
visual concreto queda delegado a la opción \texttt{estilo}.

\section{Implementación}

La implementación se basa en la clase KOMA-Script
\texttt{scrbook} \autocite{kohm2004koma}, con soporte nativo para UTF-8 mediante
LuaLaTeX.

\chapter{Conclusiones}

Se ha implementado una plantilla moderna para los trabajos de
fin de estudios de la Facultad de Informática, compatible con
LuaLaTeX y con una separación clara entre los aspectos
normativos y los estéticos.

% Normativa (TFG, V.4): se requieren conclusiones en inglés.
\chapter{Conclusions}

A modern template for the final-year projects of the Faculty of
Computer Science has been implemented, compatible with Lua\LaTeX{}
and with a clear separation between regulatory and aesthetic concerns.


\backmatter

% Normativa (TFG, V.5): en trabajos en grupo cada autor debe incluir al menos
% dos páginas describiendo su contribución personal al proyecto.
\begin{contribucion}{Frodo Bolsón}
\begin{itemize}
\item Aceptar la carga del Anillo
\item Viaje a Rivendel
\item Liderar la Comunidad del Anillo
\item Continuar hacia Mordor en solitario
\item Navegar por terrenios hostiles
\item Resistir la influencia del Anillo
\item Alcanzar la Montaña de la Perdición
\item La Purga de la Comarca
\end{itemize}
\end{contribucion}

\begin{contribucion}{Samsagaz Gamyi}
\begin{itemize}
\item Unirse a la misión de Frodo en la Comarca
\item Cuidar de Frodo durante el viaje
\item Rescatar a Frodo de la telaraña de Ella-Laraña
\item Salvar a Frodo de la Torre de Cirith Ungol
\item Ayudar a Frodo a escalar la Montaña de la Perdición
\item Participar en la Purga de la Comarca
\end{itemize}
\end{contribucion}

\printbibliography[heading=bibintoc]

\end{document}
